% =====================================================================
\chapter{Le backend en profondeur}
\label{ch:backend}

Ce chapitre traite les fichiers backend importants selon la grille demandée :
objectif, imports, classes, fonctions, paramètres, valeurs de retour, appels
internes, accès base de données, endpoints, gestion des erreurs,
authentification et dépendances. Les parties complexes sont décomposées étape
par étape.

% ---------------------------------------------------------------------
\section{Configuration --- \file{app/core/config.py}}

\subsection*{Objectif}
Fournir un objet \fn{settings} unique, validé, construit à partir des variables
d'environnement et du fichier \file{.env}.

\subsection*{Imports importants}
\texttt{pydantic\_settings.BaseSettings} et \texttt{SettingsConfigDict}.

\subsection*{Classe \cls{Settings}}

\begin{longtable}{@{}p{4.4cm}p{2.4cm}p{7.9cm}@{}}
\toprule
\textbf{Champ} & \textbf{Défaut} & \textbf{Rôle}\\
\midrule
\endhead
\texttt{database\_url} & \emph{requis} & Chaîne DSN PostgreSQL\\
\texttt{jwt\_secret\_key} & \emph{requis} & Clé de signature HS256\\
\texttt{jwt\_algorithm} & \texttt{HS256} & Algorithme JWT\\
\texttt{jwt\_expire\_minutes} & \texttt{60} & Durée de vie du jeton\\
\texttt{groq\_api\_key} & \emph{requis} & Clé de l'API Groq\\
\texttt{anthropic\_api\_key} & \texttt{""} & Déclarée, non utilisée dans le chemin actif\\
\texttt{chroma\_host} / \texttt{chroma\_port} & \texttt{localhost} / \texttt{8001} & Base vectorielle\\
\texttt{storage\_vm\_host} & \texttt{None} & Hôte SSH de la VM Storage\\
\texttt{storage\_vm\_port} & \texttt{22} & Port SSH\\
\texttt{storage\_vm\_user} & \texttt{None} & Compte SSH\\
\texttt{storage\_vm\_ssh\_key\_path} & \texttt{None} & Chemin de la clé privée\\
\texttt{upload\_dir} & \texttt{uploads} & Racine des fichiers déposés\\
\texttt{embedding\_model} & \texttt{intfloat/multilingual-e5-small} & Modèle d'embedding\\
\texttt{embedding\_query\_prefix} & \texttt{"query: "} & Préfixe E5 pour les questions\\
\texttt{embedding\_passage\_prefix} & \texttt{"passage: "} & Préfixe E5 pour les documents\\
\bottomrule
\end{longtable}

\begin{keybox}
\textbf{Décision de conception explicite dans le code.} Les quatre variables
SSH sont \emph{optionnelles}. Le commentaire l'assume : « absents,
\file{capacity.py} renvoie une erreur explicite au moment de l'appel plutôt que
de faire planter tout le backend au démarrage pour qui n'a pas encore la VM
configurée ».
\end{keybox}

% ---------------------------------------------------------------------
\section{Sécurité --- \file{app/core/security.py}}

\subsection*{Fonctions}

\begin{longtable}{@{}p{4.6cm}p{4.2cm}p{5.9cm}@{}}
\toprule
\textbf{Fonction} & \textbf{Paramètres} & \textbf{Retour et comportement}\\
\midrule
\endhead
\fn{hash\_password} & \texttt{plain\_password: str} & \texttt{str} --- empreinte bcrypt\\
\fn{verify\_password} & \texttt{plain}, \texttt{hashed} & \texttt{bool}\\
\fn{create\_access\_token} & \texttt{subject: str}, \texttt{role: str} & \texttt{str} --- JWT contenant \texttt{sub}, \texttt{role}, \texttt{exp}\\
\fn{decode\_access\_token} & \texttt{token: str} & \texttt{dict} ou \texttt{None} si \cls{JWTError} (jeton invalide \emph{ou} expiré)\\
\fn{create\_cross\_domain\_token} & \texttt{subject}, \texttt{workspace}, \texttt{minutes=30} & \texttt{str} --- JWT avec \texttt{scope="cross\_domain\_read\_only"} et \texttt{workspace}\\
\bottomrule
\end{longtable}

\subsection*{Gestion des erreurs}
\fn{decode\_access\_token} ne lève jamais : elle attrape \cls{JWTError} et
renvoie \texttt{None}. C'est l'appelant (\fn{get\_current\_user}) qui transforme
ce \texttt{None} en HTTP 401. Ce choix évite de disperser la gestion
d'exception JWT dans toutes les routes.

% ---------------------------------------------------------------------
\section{Dépendances FastAPI --- \file{app/api/deps.py}}

\subsection*{\fn{get\_current\_user} : décomposition pas à pas}

\begin{enumerate}
  \item \texttt{HTTPBearer()} extrait l'en-tête \texttt{Authorization: Bearer
  <token>}. En son absence, FastAPI répond 403 avant même d'entrer dans la
  fonction.
  \item \fn{decode\_access\_token(token)} : si le résultat est \texttt{None},
  levée d'une \cls{HTTPException} 401 avec
  \texttt{WWW-Authenticate: Bearer} et le message
  « Token invalide ou expire ».
  \item \texttt{payload.get("sub")} donne l'e-mail. \textbf{Accès base :}
  \texttt{db.query(User).filter(User.email == email).first()}.
  \item Si l'utilisateur est absent \emph{ou} \texttt{is\_active} est faux :
  401 « Utilisateur introuvable ou desactive ».
  \item Retour : l'entité \cls{User} rattachée à la session de la requête, donc
  \texttt{current\_user.role.name} est chargeable par relation.
\end{enumerate}

\begin{infobox}
\textbf{Détail important pour la soutenance.} Le rôle est présent dans le jeton,
mais \fn{require\_role} ne le lit \emph{pas} dans le jeton : il lit
\texttt{current\_user.role.name}, donc la base. Conséquence : changer le rôle
d'un utilisateur en base prend effet immédiatement, sans attendre l'expiration
de son jeton.
\end{infobox}

\subsection*{\fn{require\_role} : une fabrique de dépendances}

\begin{lstlisting}[language=Python,caption={api/deps.py --- contrôle de rôle}]
def require_role(*allowed_roles: str):
    def role_checker(current_user: User = Depends(get_current_user)) -> User:
        if current_user.role.name not in allowed_roles:
            raise HTTPException(status_code=403, detail="Acces refuse pour ce role")
        return current_user
    return role_checker
\end{lstlisting}

\noindent
Chaque routeur crée sa dépendance une seule fois, au niveau du module :
\texttt{\_storage\_access = require\_role("storage\_engineer", "administrator")}.

% ---------------------------------------------------------------------
\section{Authentification --- \file{api/v1/auth.py} et \file{services/auth\_service.py}}

\subsection*{Endpoints}

\begin{longtable}{@{}p{3.4cm}p{2.4cm}p{4.2cm}p{4.5cm}@{}}
\toprule
\textbf{Endpoint} & \textbf{Auth} & \textbf{Entrée} & \textbf{Sortie}\\
\midrule
\endhead
\route{POST /auth/login} & aucune & \cls{LoginRequest} (\texttt{EmailStr}, \texttt{password}) & \cls{TokenResponse} (\texttt{access\_token}, \texttt{token\_type="bearer"})\\
\route{GET /auth/me} & \fn{get\_current\_user} & --- & \cls{UserResponse} (\texttt{id}, \texttt{email}, \texttt{full\_name}, \texttt{role})\\
\bottomrule
\end{longtable}

\subsection*{Chaîne d'appels}
\route{POST /auth/login} \arrowto{} \fn{auth\_service.login(db, email, password)}
\arrowto{} \fn{authenticate\_user} \arrowto{} requête SQL sur \tbl{users}
\arrowto{} \fn{verify\_password} \arrowto{} \fn{create\_access\_token}.

\subsection*{Gestion des erreurs}
\fn{login} renvoie \texttt{None} dans trois cas indistincts : e-mail inconnu,
mot de passe faux, compte désactivé. La route les traduit tous en un seul 401
« Email ou mot de passe incorrect ». C'est délibéré : distinguer les cas
révélerait quels comptes existent.

\subsection*{Note sur \cls{UserResponse}}
\file{admin\_routes.py} construit ce schéma \textbf{à la main} et le commente :
« \texttt{UserResponse.role} attend une chaîne, alors que \texttt{User.role} est
une relation SQLAlchemy (objet \cls{Role}) ». Le mapping automatique
\texttt{from\_attributes} échouerait ici.

% ---------------------------------------------------------------------
\section{Le chat --- \file{api/v1/storage.py}}

\subsection*{Endpoints}

\begin{longtable}{@{}p{6.0cm}p{8.8cm}@{}}
\toprule
\textbf{Endpoint} & \textbf{Comportement}\\
\midrule
\endhead
\route{GET /storage/chat/conversations} & Fils de l'utilisateur, le plus récemment utilisé en tête\\
\route{GET /storage/chat/conversations/\{id\}/messages} & Messages du fil, avec leurs \cls{ActionRequestOut}\\
\route{PATCH /storage/chat/conversations/\{id\}} & Renomme (sans toucher à \texttt{updated\_at})\\
\route{DELETE /storage/chat/conversations/\{id\}} & 204 ; les messages suivent en cascade\\
\route{POST /storage/chat} & Un tour de dialogue complet\\
\bottomrule
\end{longtable}

\subsection*{Isolation par utilisateur}
Toutes ces routes passent par \fn{\_require\_conversation}, qui appelle
\fn{chat\_service.get\_conversation(db, user.id, conversation\_id)} --- filtré
sur \texttt{user\_id}. Le commentaire justifie le code de retour : « Le fil
demandé, ou 404 --- y compris s'il appartient à un collègue : dire "interdit"
plutôt qu'"introuvable" confirmerait son existence. »

\subsection*{\route{POST /storage/chat} : sept étapes}

\begin{enumerate}
  \item \textbf{Validation} --- \cls{ChatRequest} : \texttt{message: str},
  \texttt{conversation\_id: int | None}.
  \item \textbf{Résolution du fil} --- si \texttt{conversation\_id} est fourni,
  \fn{\_require\_conversation} puis \fn{chat\_service.history\_for\_llm} (10
  derniers messages).
  \item \textbf{Contexte d'incident} --- si le fil porte un
  \texttt{incident\_ticket\_id}, on recharge le ticket et on construit un
  \cls{IncidentContext} via \fn{incident\_investigation.context\_for}.
  \item \textbf{Appel de l'agent} ---
  \fn{await storage\_persona.handle\_message(message, created\_by=user.id,
  history=..., incident=...)}.
  \item \textbf{Création différée du fil} --- \emph{après} la réponse
  seulement. Commentaire : « Rien n'est écrit avant la réponse : si l'agent
  échoue, l'exception remonte sans laisser dans la liste un fil vide. »
  \item \textbf{Écriture du tour} --- \fn{chat\_service.save\_turn} écrit les
  deux messages dans la même transaction, puis \fn{attach\_actions} relie les
  \tbl{action\_requests} au message de l'assistant.
  \item \textbf{Réponse} --- \cls{ChatTurnOut} : \texttt{response},
  \texttt{conversation}, \texttt{user\_message}, \texttt{assistant\_message}.
\end{enumerate}

% ---------------------------------------------------------------------
\section{Le cœur de l'IA --- \file{app/personas/storage/agent.py}}

C'est le fichier le plus dense du projet (445 lignes). Il est décomposé
intégralement.

\subsection{Structures de données}

\subsubsection*{\cls{IncidentContext} (dataclass)}
Champs : \texttt{ticket\_id}, \texttt{ticket\_number}, \texttt{title},
\texttt{description}, \texttt{severity}, \texttt{affected\_system},
\texttt{mcp\_data}.

\begin{keybox}
\textbf{Pourquoi une dataclass et non l'entité ORM ?} Le commentaire l'explique :
« l'agent ouvre ses propres sessions et vit plus longtemps que celle de la
requête HTTP --- une entité détachée y deviendrait inutilisable à la première
lecture d'attribut ». C'est un problème classique de SQLAlchemy
(\cls{DetachedInstanceError}), résolu par la copie des champs.
\end{keybox}

\noindent
\textbf{Propriété \fn{technical\_id}} --- cherche successivement
\texttt{job\_id}, \texttt{volume\_id}, \texttt{mountpoint} dans
\texttt{mcp\_data}. Elle distingue l'\emph{identifiant technique}
(\texttt{backup-app-critical.timer}) du \emph{libellé affiché}
(\texttt{Sauvegarde de l'application critique}, qui est la \texttt{Description}
systemd). Confondre les deux ferait échouer l'action au moment précis où
l'ingénieur la confirme.

\noindent
\textbf{Méthode \fn{as\_prompt\_block}} --- produit le bloc de prompt injecté en
mode investigation : titre, ressource, cible technique, constat initial, relevé
MCP en JSON, puis trois consignes numérotées (lire les preuves, énoncer la cause
racine, proposer \emph{une} remédiation en appelant réellement l'outil).

\subsubsection*{\cls{AgentTurn} (dataclass)}
\texttt{text: str} et \texttt{proposed\_action\_ids: list[uuid.UUID]}. Le
commentaire dit pourquoi ce type existe : « \fn{handle\_message} renvoyait une
simple chaîne. Les propositions d'action doivent remonter jusqu'à la route pour
être rattachées au message qui les porte --- sans quoi l'interface ne saurait
pas quels boutons afficher. »

\subsection{\fn{\_propose\_action} --- le garde-fou central}

\textbf{Paramètres} : \texttt{db}, \texttt{tool\_name}, \texttt{arguments},
\texttt{created\_by}, \texttt{incident}.\\
\textbf{Retour} : \texttt{tuple[dict, uuid.UUID | None]} --- le dictionnaire est
réinjecté dans la conversation comme « résultat » de l'outil, l'UUID sert à
rattacher la proposition au message.

\begin{enumerate}
  \item \textbf{Extraction de la cible} parmi \texttt{("target", "volume\_id",
  "job\_id")}. Absente : retour \texttt{\{"status": "error", ...\}} et pas
  d'action créée.
  \item \textbf{Correction de la cible} : si \texttt{target ==
  incident.affected\_system} et qu'un \fn{technical\_id} existe, on substitue
  l'identifiant technique, avec journalisation
  (\texttt{"Cible de \%s corrigee : \%r -> \%r"}).
  \item \textbf{Nettoyage} : la clé \texttt{confirm} est retirée des paramètres.
  Elle sera ajoutée par \fn{confirm\_and\_execute\_action}, et la laisser ici la
  passerait deux fois.
  \item \textbf{Écriture} : \fn{incident\_service.create\_action\_request(...)}
  \arrowto{} une ligne \tbl{action\_requests} au statut \texttt{PENDING}.
  \item \textbf{Retour au modèle} : un dictionnaire dont le message est
  volontairement directif ---
  \emph{« Action NOT executed: it has been proposed to the engineer, who
  confirms or rejects it from the interface. Present it as a recommendation,
  never as a fact. »}
\end{enumerate}

\subsection{\fn{handle\_message} --- la boucle de tool calling}

\textbf{Signature} : \fn{async handle\_message(user\_message, created\_by=None,
history=None, incident=None) -> AgentTurn}.

\begin{figure}[h!]
\centering
\begin{tikzpicture}[node distance=5mm and 9mm]
\node[bluebox,minimum width=34mm] (rag) {1. RAG\\\tiny retriever.retrieve(storage\_kb, question, n=3)};
\node[bluebox,below=of rag,minimum width=34mm] (prompt) {2. Prompt système\\\tiny STORAGE\_SYSTEM\_PROMPT + date + contexte RAG + bloc incident};
\node[bluebox,below=of prompt,minimum width=34mm] (hist) {3. Historique\\\tiny seuls les tours user/assistant sont rejoués};
\node[ambox,below=of hist,minimum width=34mm] (open) {4. \texttt{async with StorageMCPClient()}\\\tiny list\_tools\_for\_llm()};
\node[vibox,below=of open,minimum width=34mm] (loop) {5. Boucle (max 6 tours)\\\tiny groq\_client.chat\_with\_tools(messages, tools)};
\node[greenbox,right=16mm of loop,minimum width=30mm] (read) {5a. Outil de lecture\\\tiny mcp.call\_tool(...)};
\node[redbox,below=6mm of read,minimum width=30mm] (rem) {5b. Outil de remédiation\\\tiny \_propose\_action(...)};
\node[bluebox,below=16mm of loop,minimum width=34mm] (rules) {6. evaluate\_tool\_result\\\tiny détection d'anomalie};
\node[bluebox,below=of rules,minimum width=34mm] (recover) {7. Rattrapage\\\tiny \_recover\_written\_call()};
\node[bluebox,below=of recover,minimum width=34mm] (tickets) {8. Création des tickets\\\tiny create\_ticket\_from\_diagnosis};
\node[box,below=of tickets,minimum width=34mm,draw=oasisdark] (out) {\textbf{AgentTurn}\\\tiny text + proposed\_action\_ids};

\draw[flowb] (rag)--(prompt); \draw[flowb] (prompt)--(hist);
\draw[flowb] (hist)--(open);  \draw[flowb] (open)--(loop);
\draw[flow] (loop)--(read);   \draw[flow] (loop)-|(rem);
\draw[flow] (read.east) to[bend right=45] node[lbl,right]{résultat} (rules.east);
\draw[flow] (rem.east) to[bend right=35] (rules.east);
\draw[flowb] (loop)--(rules); \draw[flowb] (rules)--(recover);
\draw[flowb] (recover)--(tickets); \draw[flowb] (tickets)--(out);
\end{tikzpicture}
\caption{Les huit temps de \fn{handle\_message}.}
\end{figure}

\clearpage

\subsubsection*{Étape 1 --- RAG}
\fn{retriever.retrieve(collection\_name="storage\_kb", question=user\_message,
n\_results=3)} \arrowto{} \fn{get\_context\_for\_query} \arrowto{}
\fn{vectorstore.search} \arrowto{} \fn{embed\_query} (préfixe
\texttt{"query: "}) \arrowto{} \texttt{collection.query}. Le résultat est une
chaîne où chaque passage est préfixé de \texttt{[Source: <fichier>]}.

\subsubsection*{Étape 2 --- Prompt système}
Il concatène quatre blocs : le \fn{STORAGE\_SYSTEM\_PROMPT}, la date et l'heure
du serveur, une consigne explicite sur MCP, et le contexte RAG. En mode
investigation s'ajoute \fn{incident.as\_prompt\_block()}.

\begin{lstlisting}[caption={Consigne anti-hallucination injectée à chaque tour}]
Use MCP only for current storage information. The tools that change the VM
(run_backup, create_snapshot, restore_from_backup) never execute from the chat:
calling one records an action proposal that the engineer confirms with a click.
So call them normally when you want to propose one, but never copy a tool call
as JSON into your answer -- written that way it creates nothing. And never say
a remediation has been launched: it is waiting on a human decision.
\end{lstlisting}

\subsubsection*{Étape 3 --- Historique}
Seuls les messages de rôle \texttt{user} et \texttt{assistant} sont rejoués. Le
commentaire justifie : « les appels d'outils d'une conversation précédente ne
sont pas rejouables (leurs résultats ne sont plus à jour), le modèle les
redemande s'il en a besoin ».

\subsubsection*{Étape 4 --- Session MCP et session d'audit}
\texttt{async with StorageMCPClient() as mcp} démarre le sous-processus.
\texttt{audit\_db = SessionLocal() if created\_by is not None else None} : sans
utilisateur identifié, on n'audite pas plutôt que d'inventer un auteur.

\subsubsection*{Étape 5 --- La boucle (au plus \fn{\_MAX\_TOOL\_ROUNDS} = 6)}

\begin{enumerate}
  \item \fn{groq\_client.chat\_with\_tools(messages, tools)}.
  \item Si \texttt{reply.tool\_calls} est vide \arrowto{} \texttt{final\_text =
  reply.content}, sortie de boucle.
  \item Sinon, le message assistant \emph{avec ses \texttt{tool\_calls}} est
  réinjecté dans \texttt{messages} (indispensable au protocole OpenAI/Groq).
  \item Pour chaque appel : \texttt{json.loads(arguments)} dans un
  \texttt{try/except JSONDecodeError} qui retombe sur \texttt{\{\}}.
  \item Aiguillage : \texttt{REMEDIATION\_TOOLS} \arrowto{} proposition ;
  sinon \arrowto{} exécution MCP réelle, avec \texttt{RuntimeError} traduite en
  \texttt{\{"status": "error", "message": ...\}} plutôt que propagée.
  \item Le résultat est réinjecté sous forme d'un message de rôle
  \texttt{tool}, avec le \texttt{tool\_call\_id} correspondant.
\end{enumerate}

\noindent
Si la boucle atteint 6 tours sans réponse finale, un dernier appel sans outils
produit le texte, avec un message de repli explicite (« MCP collection limit
reached; try a more targeted question »).

\subsubsection*{Étape 6 --- Détection déterministe}
\fn{evaluate\_tool\_result(name, result)} applique les règles de
\file{incident\_rules.py}. \textbf{Aucune décision d'incident n'est prise par le
modèle.}

\subsubsection*{Étape 7 --- \fn{\_recover\_written\_call} : le rattrapage}

C'est une des idées les plus fines du projet. Le défaut visé est décrit dans le
code : « le diagnostic est juste, la remédiation est la bonne, mais elle est
recopiée en JSON dans la réponse. Écrite ainsi elle ne crée rien, et l'ingénieur
lit une analyse complète sans avoir aucun bouton à cliquer ».

\begin{enumerate}
  \item Une expression régulière \fn{\_JSON\_OBJECT} extrait les objets JSON du
  texte (un niveau d'imbrication toléré).
  \item \fn{\_candidate\_calls} accepte trois formes rencontrées selon les
  modèles : \texttt{\{"tool", "args"\}}, \texttt{\{"name", "arguments"\}},
  \texttt{\{"tool", "parameters"\}}, plus la forme où le nom de l'outil est
  directement une clé.
  \item \textbf{Double validation} : l'outil doit appartenir à
  \texttt{REMEDIATION\_TOOLS}, \emph{et} la cible doit être
  \texttt{incident.affected\_system} ou \texttt{incident.technical\_id}. Une
  cible inventée ou portant sur une autre machine est ignorée.
  \item Ce qui en sort reste une proposition \texttt{pending}, soumise au même
  clic humain.
\end{enumerate}

\subsubsection*{Étape 8 --- Création des tickets détectés}
Une nouvelle session (\texttt{with SessionLocal() as db}) écrit les tickets via
\fn{create\_ticket\_from\_diagnosis}. Sans \texttt{created\_by}, on journalise un
avertissement et on ne crée rien.

\begin{warnbox}
\textbf{Gestion d'erreurs à souligner.} Chaque bloc à risque est entouré d'un
\texttt{try/except} qui journalise via \fn{logger.exception} et \textbf{ne fait
jamais échouer le tour}. Une panne d'audit, de création de ticket ou de
rattrapage ne doit pas priver l'ingénieur de sa réponse.
\end{warnbox}

% ---------------------------------------------------------------------
\section{Client LLM --- \file{app/llm/groq\_client.py}}

\subsection*{Classe \cls{GroqClient}}
\texttt{self.client = Groq(api\_key=settings.groq\_api\_key)} ;
\texttt{self.model = getattr(settings, "groq\_model", "openai/gpt-oss-120b")}.

\begin{infobox}
\texttt{groq\_model} n'est \textbf{pas} un champ de \cls{Settings} : le
\fn{getattr} retombe donc toujours sur le défaut
\texttt{"openai/gpt-oss-120b"}. C'est une lecture du code, pas une supposition.
\end{infobox}

\subsection*{\fn{chat\_with\_tools(messages, tools=None)}}
Construit une liste de modèles candidats
(\texttt{openai/gpt-oss-120b}, \texttt{openai/gpt-oss-20b},
\texttt{llama-3.3-70b-versatile}), dédoublonnée en conservant l'ordre, puis les
essaie successivement. À la première réussite, \texttt{self.model} est mis à
jour --- le modèle qui fonctionne devient le défaut pour les appels suivants. Si
tous échouent, la dernière exception est relevée.

\noindent
Options envoyées : \texttt{temperature=0.3} et, si des outils sont fournis,
\texttt{tools} et \texttt{tool\_choice="auto"}.

\subsection*{\fn{chat(prompt, system\_prompt=None)}}
Enveloppe simple qui renvoie \texttt{.content}. Utilisée par
\fn{incident\_investigation.\_summarise} pour la synthèse
\texttt{root\_cause} / \texttt{impact\_summary}.

% ---------------------------------------------------------------------
\section{Actions sensibles --- \file{api/v1/actions.py}}

\subsection*{Endpoints}

\begin{longtable}{@{}p{4.6cm}p{2.6cm}p{7.5cm}@{}}
\toprule
\textbf{Endpoint} & \textbf{Auth} & \textbf{Comportement}\\
\midrule
\endhead
\route{GET /storage/actions} & \fn{get\_current\_user} & Actions de l'utilisateur courant (filtre \texttt{requested\_by}), pagination \texttt{limit} (\(\le 200\)) et \texttt{offset}\\
\route{POST /storage/actions/} & \fn{get\_current\_user} & Enregistre une proposition (\texttt{pending}). \cls{ValueError} \arrowto{} 400\\
\route{POST /storage/actions/\{id\}/confirm} & \fn{get\_current\_user} & \textbf{Seul chemin d'exécution réelle}\\
\route{POST /storage/actions/\{id\}/reject} & \fn{get\_current\_user} & Passe en \texttt{REJECTED} et écrit une note sur le ticket\\
\route{GET /storage/actions/\{id\}} & \fn{get\_current\_user} & Lecture d'une action\\
\bottomrule
\end{longtable}

\subsection*{\fn{confirm\_action} pas à pas}

\begin{enumerate}
  \item \texttt{mcp = StorageMCPClient()} puis \texttt{await mcp.start()}. Si le
  sous-processus ne démarre pas : \textbf{HTTP 503}, et l'action \emph{reste}
  \texttt{PENDING}, donc rejouable.
  \item \fn{await incident\_service.confirm\_and\_execute\_action(db,
  action\_id, user\_id, mcp\_client=mcp)}.
  \item \fn{await \_close\_the\_loop(db, action, user, mcp)}.
  \item \texttt{finally: await mcp.stop()} --- le sous-processus est toujours
  fermé.
\end{enumerate}

\begin{keybox}
\textbf{Le correctif documenté dans le fichier.} Le docstring raconte un bug
réel : « cette route utilisait une instance module-level de
\cls{StorageMCPClient} dont personne n'appelait jamais \fn{start()} \ldots{}
\fn{call\_tool()} levait donc systématiquement "Storage MCP indisponible",
l'action retombait en FAILED et l'API répondait quand même 200 : le seul chemin
autorisé pour exécuter une action destructive ne fonctionnait pas. »
C'est un excellent exemple à citer en soutenance.
\end{keybox}

\subsection*{\fn{\_close\_the\_loop} : la fermeture du cycle}

\begin{enumerate}
  \item Sans \texttt{incident\_ticket\_id} \arrowto{} rien à faire.
  \item Si l'action est \texttt{FAILED} \arrowto{} note
  « \ldots{} confirmed by X but failed: <raison> », et le ticket reste ouvert.
  \item Sinon \arrowto{} \fn{await incident\_verification.verify(ticket,
  mcp=mcp, db=db, user\_id=user.id)}.
  \item Une note consigne le verdict ; si \texttt{verdict.resolved} est vrai,
  \fn{incident\_service.resolve\_ticket}.
  \item \texttt{except Exception} global : la fermeture du cycle ne doit
  \textbf{jamais} faire échouer la requête. « L'action, elle, a bien eu lieu, et
  la surveillance périodique repassera vérifier. »
\end{enumerate}

% ---------------------------------------------------------------------
\section{Service transactionnel --- \file{services/incident\_service.py}}

\subsection*{\fn{confirm\_and\_execute\_action} : le point critique}

\begin{lstlisting}[language=Python,caption={Extrait réel --- transition, exécution, qualification du résultat}]
if action.status != ActionRequestStatus.PENDING:
    raise ValueError(f"Action already handled (current status: {action.status.value})")

action.status = ActionRequestStatus.CONFIRMED
action.confirmed_by = user_id
action.confirmed_at = datetime.utcnow()
db.commit()

try:
    result = await mcp_client.call_tool(
        action.action_type, db=db, user_id=user_id,
        action_request_id=action.id, confirm=True, **action.parameters,
    )
    if isinstance(result, dict) and result.get("error"):
        action.status = ActionRequestStatus.FAILED
    else:
        action.status = ActionRequestStatus.COMPLETED
    action.result = result
except Exception as error:
    action.status = ActionRequestStatus.FAILED
    action.result = {"error": str(error)}
\end{lstlisting}

\begin{keybox}
\textbf{Trois protections dans dix lignes.}

\emph{(1)} Le test \texttt{status != PENDING} empêche la double exécution : deux
clics sur « Confirm » ne lancent pas deux restaurations.

\emph{(2)} Le passage en \texttt{CONFIRMED} est \texttt{commit} \textbf{avant}
l'appel, donc la trace de l'intention survit à un crash pendant l'exécution.

\emph{(3)} Le test \texttt{result.get("error")} est essentiel : les outils
Storage ne lèvent pas d'exception quand la VM refuse une commande, ils renvoient
\texttt{\{"error": ...\}}. Sans ce test, \texttt{systemctl start} sur un job qui
échoue aussitôt marquerait l'action \texttt{COMPLETED}, et
\fn{has\_completed\_action} prendrait une remédiation ratée pour une remédiation
faite --- jusqu'à clore le ticket sur cette base.
\end{keybox}

\subsection*{\fn{\_generate\_ticket\_number}}
Le numéro suivant est dérivé du \textbf{maximum} existant, pas d'un
\texttt{COUNT}. Justification donnée : « compter suppose une suite sans trou. Un
seul ticket supprimé, et le compteur repointe sur un numéro déjà pris ---
l'insertion viole alors l'unicité de \texttt{ticket\_number} et fait tomber tout
le passage de surveillance, donc la détection elle-même. » Format :
\texttt{INC-<année>-0001}.

\subsection*{\fn{\_open\_ticket\_for} : déduplication}
Recherche d'abord par identifiant technique dans le JSONB
(\texttt{mcp\_data["job\_id"]}, puis \texttt{mcp\_data["volume\_id"]}), et
\emph{seulement ensuite} par \texttt{affected\_system}. Motif : renommer l'unit
sur la VM suffisait à faire ouvrir un second ticket pour le même job.

\subsection*{\fn{has\_completed\_action}}
Retourne vrai s'il existe au moins une \tbl{action\_requests} en
\texttt{COMPLETED} sur ce ticket. C'est la garde de l'auto-résolution : « un
incident ne se referme pas tout seul parce que la VM va mieux, mais parce
qu'une action confirmée par un humain a produit cet effet ».

% ---------------------------------------------------------------------
\section{Investigation --- \file{services/incident\_investigation.py}}

\subsection*{\fn{start(db, ticket, user)} --- séquence}

\begin{enumerate}
  \item \fn{find\_conversation(db, ticket\_id, user\_id)} : un fil existant est
  renvoyé tel quel, sans relancer l'agent. Le fil est \textbf{par utilisateur},
  car une conversation appartient à une personne et n'est lisible que par elle.
  \item \fn{context\_for(ticket)} \arrowto{} \cls{IncidentContext}.
  \item \fn{\_investigation\_prompt(ticket)} rédige la question à la place de
  l'ingénieur --- et elle est enregistrée comme un vrai message utilisateur,
  « en relisant le fil six mois plus tard, on voit la question qui a lancé
  l'analyse, et non un diagnostic tombé de nulle part ».
  \item L'agent tourne \textbf{avant} toute écriture.
  \item Création du fil, rattachement à \texttt{incident\_ticket\_id},
  \fn{save\_turn}.
  \item Si l'agent n'a proposé aucune action \arrowto{} \fn{\_fallback\_action}.
  \item \fn{attach\_actions}, puis \fn{mark\_in\_progress} (le ticket passe de
  \texttt{OPEN} à \texttt{IN\_PROGRESS}).
  \item \fn{\_summarise(turn.text)} : second appel Groq qui renvoie un JSON
  \texttt{\{"root\_cause", "impact\_summary"\}} recopié sur le ticket.
  \item Note d'audit nommée, puis retour \texttt{\{"conversation\_id",
  "created"\}}.
\end{enumerate}

\subsection*{\fn{\_fallback\_action} : filet de sécurité déterministe}
Parcourt \texttt{ticket.recommendations} (calculées par
\file{incident\_rules.py} à la détection), retient la première dont
\texttt{action\_type} appartient à \texttt{REMEDIATION\_TOOLS}, écarte les
doublons (\texttt{pending}, \texttt{confirmed}, \texttt{completed}) et crée
l'action.

\begin{infobox}
Le commentaire résume la philosophie : « on ne demande donc au modèle que ce que
lui seul sait faire, le diagnostic, et on tient la proposition de ce qui est
déterministe ».
\end{infobox}

\subsection*{\fn{\_parse\_summary}}
Découpe entre la première \texttt{\{} et la dernière \texttt{\}} avant
\texttt{json.loads}, pour survivre à un modèle qui aurait ajouté un bloc de code
markdown. Ne conserve que les deux clés attendues.

% ---------------------------------------------------------------------
\section{Vérification --- \file{services/incident\_verification.py}}

\subsection*{Classe \cls{Verdict}}
Trois attributs : \texttt{resolved: bool}, \texttt{detail: str},
\texttt{conclusive: bool}. Le troisième distingue « toujours en panne » de
« état pas encore lisible » (job encore en cours).

\subsection*{\fn{verify(ticket, mcp, db, user\_id)} --- aiguillage}
Le type d'incident se déduit de \texttt{ticket.mcp\_data} :
\texttt{job\_id} \arrowto{} \fn{\_verify\_backup} ;
\texttt{volume\_id} ou \texttt{mountpoint} \arrowto{} \fn{\_verify\_capacity} ;
sinon \arrowto{} verdict non concluant, clôture manuelle.

\subsection*{\fn{\_verify\_backup} --- boucle d'attente}
5 tentatives espacées de 3 secondes tant que le job est \texttt{running}, car
« un job systemd de type \texttt{oneshot} est encore \texttt{running} juste
après un \texttt{systemctl start} ». Puis la \textbf{même règle} que la
détection est réappliquée : si \fn{check\_backup\_status(job)} renvoie
\texttt{None}, le problème a disparu.

\begin{keybox}
\textbf{Point d'architecture élégant.} La vérification ne réécrit pas une
variante des seuils : elle réutilise \fn{check\_capacity} et
\fn{check\_backup\_status}. « Ce qui a ouvert le ticket est exactement ce qui
doit avoir disparu pour le fermer. »
\end{keybox}

% ---------------------------------------------------------------------
\section{Surveillance de fond --- \file{services/incident\_watcher.py}}

\subsection*{\fn{scan\_once()} --- un passage, sans exception}

\begin{enumerate}
  \item \fn{list\_backups()} --- appel \textbf{direct}, pas via MCP.
  \item Si \texttt{error} \arrowto{} rapport
  \texttt{\{"error", "created": [], "skipped": [], "closed": []\}}. Une VM
  éteinte est « une situation normale, pas une panne de l'application ».
  \item \fn{\_system\_user\_id(db)} --- compte de service.
  \item \fn{\_verify\_resolved} \textbf{d'abord} : referme les tickets dont la
  remédiation a produit son effet. L'ordre est délibéré, « un ticket clos au
  même cycle ne doit pas être rouvert par la boucle de détection qui suit ».
  \item Pour chaque job en échec : \fn{check\_backup\_status}, test
  \fn{\_already\_handled}, puis \fn{create\_ticket\_from\_diagnosis} et une
  première note d'audit.
\end{enumerate}

\subsection*{Le compte de service}
\texttt{oasis-system@oasis.local}, rôle \texttt{storage\_engineer},
\texttt{is\_active=False}, mot de passe aléatoire jamais communiqué. Il ne peut
pas servir à se connecter (\fn{get\_current\_user} rejette les comptes
désactivés) : il n'existe que pour être l'auteur des tickets automatiques, car
\texttt{incident\_tickets.created\_by} est \texttt{NOT NULL}.

\subsection*{\fn{\_already\_handled} --- anti-rouvertures}
Sans ce test, résoudre manuellement un ticket sur un job encore en échec en
rouvrirait un identique au cycle suivant, « le bouton résoudre n'aurait aucun
effet visible ». Règle : on ne rouvre qu'à partir d'une exécution
\emph{postérieure} au ticket précédent --- un ticket par échec, pas par cycle.

% ---------------------------------------------------------------------
\section{Dashboard --- \file{app/dashboard/kpi\_service.py}}

\subsection*{Objectif}
Agréger, \textbf{sans aucun appel LLM}, la santé de la VM et les compteurs de la
base. Le fichier le déclare : « Aucun appel LLM ici : le dashboard doit rester
rapide et déterministe. »

\subsection*{Le cache de santé : trois fonctions coopérantes}

\begin{longtable}{@{}p{4.6cm}p{10.2cm}@{}}
\toprule
\textbf{Fonction} & \textbf{Comportement}\\
\midrule
\endhead
\fn{refresh\_health\_snapshot} & Prend \fn{\_refresh\_lock}, revérifie la fraîcheur, appelle \fn{get\_storage\_health()}, écrit le cache sous \fn{\_health\_lock}\\
\fn{\_storage\_health(force)} & Sert le dernier relevé ; s'il dépasse 60\,s et qu'aucun rafraîchissement n'est en cours, lance un thread de fond et rend la main immédiatement\\
\fn{\_refresh\_in\_background} & Cible du thread, avale les exceptions (la VM peut être injoignable)\\
\fn{\_invalidate\_health\_snapshot} & Vide le cache, utilisé par \texttt{?refresh=true}\\
\bottomrule
\end{longtable}

\begin{infobox}
\textbf{Motif documenté.} « \fn{get\_storage\_health()} ouvre deux connexions SSH
vers la VM Storage : la réponse met plusieurs secondes, et le dashboard la
demandait à chaque affichage --- c'est ce qui rendait le changement de page si
lent. » Le verrou évite qu'« dix onglets ouverts en même temps sur un cache vide
déclenchent dix sessions SSH pour la même donnée ».
\end{infobox}

\subsection*{Constructeurs de blocs}

\begin{itemize}
  \item \fn{\_build\_capacity} --- totalise les volumes, calcule
  \texttt{reserved\_gb} (l'écart entre total et \(used + available\)
  correspondant aux blocs réservés au superutilisateur), attribue un statut par
  volume aux seuils 80\,\% / 90\,\%.
  \item \fn{\_build\_backups} --- compte
  \texttt{successful} / \texttt{failed} / \texttt{running} / \texttt{scheduled},
  et distingue \texttt{configured=False} (aucun timer) de \texttt{error} (VM
  injoignable).
  \item \fn{\_build\_incidents} --- agrégation SQL \texttt{GROUP BY severity} sur
  les statuts ouverts, plus le compte des résolus sur 7 jours.
  \item \fn{\_health\_score} --- \(100 - \sum \text{poids}\) avec
  \texttt{critical}=12, \texttt{high}=6, \texttt{medium}=2, \texttt{low}=1, plus
  une pénalité de capacité. Renvoie \texttt{None} si la VM est injoignable :
  « afficher 100/100 sans avoir rien pu mesurer donnerait une fausse
  assurance ».
  \item \fn{\_headline} --- une phrase en langage naturel, par ordre de priorité
  décroissante.
  \item \fn{get\_insights} --- observations avec priorité, détail, et lien vers
  l'écran qui permet d'agir.
\end{itemize}

% ---------------------------------------------------------------------
\section{Métriques --- \file{tools/performance.py} et \file{services/metrics\_service.py}}

\subsection*{\fn{sample\_disk\_performance()}}
\texttt{/proc/diskstats} n'expose que des compteurs cumulatifs : une seule
lecture ne dit rien. La fonction exécute donc, \textbf{en une seule session
SSH} :

\begin{lstlisting}[caption={La commande réellement envoyée à la VM}]
cat /proc/diskstats; echo '---'; sleep 2; cat /proc/diskstats
\end{lstlisting}

\noindent
Puis \fn{\_parse\_diskstats} filtre les périphériques virtuels
(\texttt{loop}, \texttt{ram}, \texttt{sr}, \texttt{fd}) et les partitions
(regex \fn{\_PARTITION}) --- « leurs E/S sont déjà comptées dans le disque
parent, les inclure reviendrait à compter deux fois ».
\fn{\_rates} calcule IOPS, débit (secteurs \(\times\) 512 octets) et latence
moyenne, et renvoie \texttt{None} si un compteur a reculé (redémarrage entre les
deux lectures).

\subsection*{\fn{read\_series(db, range\_key)} --- agrégation en deux passes}
\texttt{RANGE\_WINDOWS} définit fenêtre et granularité :
\texttt{24h}\,\arrowto{}\,1\,h, \texttt{7d}\,\arrowto{}\,6\,h,
\texttt{30d}\,\arrowto{}\,1\,jour. La largeur vient de cette table, jamais de
l'utilisateur --- ce qui ferme toute injection dans le \texttt{literal\_column}.

\begin{enumerate}
  \item Sous-requête : moyenne \emph{par device} et par intervalle.
  \item Requête principale : somme des devices pour IOPS et débit, moyenne pour
  la latence.
\end{enumerate}

\noindent
Justification : « agréger en une seule passe donnerait un total faussé dès
qu'un device a plus d'échantillons qu'un autre ».
La tendance \texttt{iops\_trend\_percent} compare la seconde moitié des points à
la première.

\begin{keybox}
\textbf{Le point d'honnêteté.} \fn{get\_performance} renvoie
\texttt{collecting = not series["points"]}. Tant que le collecteur n'a rien
enregistré, la série est vide et l'interface l'annonce, au lieu de fabriquer une
courbe.
\end{keybox}

% ---------------------------------------------------------------------
\section{Data Hub --- \file{services/knowledge\_service.py}}

\subsection*{\fn{create\_document} --- validations en cascade}

\begin{enumerate}
  \item Fichier vide \arrowto{} \cls{DocumentError}.
  \item Taille \(>\) \texttt{MAX\_UPLOAD\_BYTES} (10\,Mo) \arrowto{}
  \cls{DocumentError}.
  \item Extension hors de \texttt{ALLOWED\_EXTENSIONS}
  (\texttt{.md .txt .log .json .yaml .yml .csv .pdf}) \arrowto{}
  \cls{DocumentError}.
  \item Écriture sur disque sous \texttt{<uuid><extension>} --- \textbf{jamais}
  sous le nom fourni par le client. C'est ce qui rend toute traversée de
  répertoire impossible.
  \item Ligne \tbl{knowledge\_documents} en statut \texttt{PENDING}, commit.
  \item Extraction du texte, \fn{\_index\_document}, statut \texttt{INDEXED} et
  \texttt{indexed\_at}.
  \item En cas d'échec : statut \texttt{FAILED} et message tronqué à 500
  caractères dans \texttt{error}. \textbf{Le document reste en base}, l'interface
  affiche pourquoi il n'est pas indexé.
\end{enumerate}

\subsection*{\fn{\_index\_document}}
Refuse un document de moins de \texttt{MIN\_INDEXABLE\_WORDS} (20) mots utiles,
avec un message qui diagnostique le cas courant : « S'il s'agit d'un PDF scanné
ou d'une image exportée en PDF, il n'a pas de couche texte : il faut lui
appliquer un OCR avant de le déposer. » Puis \fn{chunk\_text} (500 mots,
recouvrement 50), \fn{embed\_passages}, et \texttt{collection.upsert} avec des
identifiants dérivés de l'UUID du document --- ce qui rend la suppression et la
réindexation possibles.

% ---------------------------------------------------------------------
\section{Demandes de service --- \file{services/request\_service.py}}

\subsection*{Les deux règles qui portent toute la logique}

\begin{enumerate}
  \item \textbf{Le sens est déduit du rôle de l'auteur, jamais fourni par le
  client.} Un \texttt{storage\_engineer} ne peut créer qu'une
  \texttt{ENGINEER\_TO\_ADMIN} ; un \texttt{administrator}, qu'une
  \texttt{ADMIN\_TO\_ENGINEER}. « Le client ne peut donc pas fabriquer une
  demande qui semblerait venir de l'autre camp. »
  \item \textbf{Seul le destinataire fait avancer le statut.}
  \texttt{request.assigned\_to != user.id} \arrowto{} \cls{RequestForbidden}.
\end{enumerate}

\subsection*{Table \fn{ALLOWED\_TRANSITIONS}}

\begin{longtable}{@{}p{3.0cm}p{5.4cm}p{6.2cm}@{}}
\toprule
\textbf{Statut} & \textbf{Ingénieur \arrowto{} Admin} & \textbf{Admin \arrowto{} Ingénieur}\\
\midrule
\endhead
\texttt{pending} & \texttt{in\_review}, \texttt{approved}, \texttt{rejected} & \texttt{acknowledged}, \texttt{in\_progress}\\
\texttt{in\_review} & \texttt{approved}, \texttt{rejected} & ---\\
\texttt{approved} & \texttt{completed} & ---\\
\texttt{acknowledged} & --- & \texttt{in\_progress}, \texttt{completed}\\
\texttt{in\_progress} & --- & \texttt{completed}\\
\texttt{rejected} & \(\varnothing\) (terminal) & ---\\
\texttt{completed} & \(\varnothing\) (terminal) & \(\varnothing\) (terminal)\\
\bottomrule
\end{longtable}

\noindent
Une transition refusée produit un message qui \emph{énumère} les transitions
possibles, ce qui rend l'erreur exploitable côté interface.

\subsection*{Horodatages distincts}
\texttt{decided\_at} est figé à l'arrivée dans \texttt{APPROVED} ou
\texttt{REJECTED}, distinct de \texttt{updated\_at} : « le détail annonce
"refusée le \ldots", pas "modifiée le \ldots" ».

\subsection*{Pièces jointes}
Plafond 10\,Mo, éventail d'extensions plus large que le Data Hub (« ce qui
circule ici n'a pas vocation à être indexé, seulement lu par l'autre partie »),
stockage sous UUID dans \file{upload\_dir/requests}, suppression réservée au
déposant.

% ---------------------------------------------------------------------
\section{Recherche globale --- \file{services/search\_service.py}}

\subsection*{\fn{search(db, user, query, limit\_per\_type)}}
Refuse les requêtes de moins de 2 caractères (« elle renverrait la moitié de la
base pour rien »), puis concatène quatre recherches :
incidents, demandes, conversations, documents.

\subsection*{\fn{\_pattern} --- échappement}

\begin{lstlisting}[language=Python,caption={Protection contre les jokers SQL}]
escaped = query.replace("\\", "\\\\").replace("%", "\\%").replace("_", "\\_")
return f"%{escaped}%"
\end{lstlisting}

\noindent
Puis \texttt{.ilike(pattern, escape="\textbackslash\textbackslash")}. Sans cela,
« un \texttt{\%} tapé dans la barre de recherche deviendrait un joker et
remonterait toute la base ».

\subsection*{Cloisonnement}
Les demandes et les conversations sont filtrées sur l'utilisateur courant
(\texttt{created\_by == user\_id or assigned\_to == user\_id} pour les demandes,
\texttt{user\_id} pour les conversations). Les incidents et les documents sont
communs.

\subsection*{Identifiants retournés}
Pour un incident, c'est l'\textbf{UUID} qui est renvoyé et non le numéro
affiché : « la page Incidents recharge le ticket par la route
\route{/incidents/\{id\}}, ce qui l'ouvre même quand il ne figure pas dans la
liste actuellement filtrée ».

% ---------------------------------------------------------------------
\section{Rapports PDF --- \file{services/report\_service.py}}

\subsection*{Deux entrées}
\fn{build\_global\_report(db, author)} et
\fn{build\_incident\_report(db, ticket\_id, author)} renvoient des
\texttt{bytes}. La seconde renvoie \texttt{None} si le ticket n'existe pas,
traduit en 404 par la route.

\subsection*{Mise en forme}
Le fichier définit une charte : couleurs (\texttt{\_BRAND}, \texttt{\_INK},
\texttt{\_MUTED}, \texttt{\_LINE}, \texttt{\_ZEBRA}), styles de paragraphe
(\fn{\_styles}), un constructeur de tableau (\fn{\_table}), un pied de page
(\fn{\_footer}) et un moteur de rendu (\fn{\_render}).

\subsection*{Détail HTTP à connaître}
La route ajoute \texttt{Access-Control-Expose-Headers: Content-Disposition}.
Sans cet en-tête, « le nom de fichier n'est pas lisible par le frontend, qui lit
la réponse en blob depuis un fetch authentifié ».

% ---------------------------------------------------------------------
\section{Audit HTTP --- \file{middleware/audit\_middleware.py}}

\subsection*{\fn{dispatch}}

\begin{enumerate}
  \item Court-circuit sur \texttt{EXCLUDED\_PATHS} (\file{/docs},
  \file{/openapi.json}, \file{/redoc}, \file{/favicon.ico}).
  \item Chronomètre \texttt{time.monotonic()} autour de \fn{call\_next}.
  \item \textbf{Session dédiée} \fn{SessionLocal()} --- pas celle de la requête,
  qui est déjà refermée à ce stade.
  \item \fn{\_extract\_user\_id} redécode le \texttt{Bearer} pour attribuer la
  ligne.
  \item Écriture d'une \cls{AuditLog} : \texttt{action} = méthode,
  \texttt{resource} = chemin, \texttt{details} =
  \texttt{\{status\_code, duration\_ms, client\_ip\}}.
  \item \texttt{except Exception: db.rollback()} --- un échec d'audit ne casse
  jamais la réponse.
\end{enumerate}

\begin{keybox}
\textbf{Le corps n'est jamais journalisé.} Commentaire du fichier : « on ne
journalise jamais le corps de la requête (donc jamais un mot de passe en clair
envoyé à \route{/auth/login}) --- seulement des métadonnées ». En contrepartie,
les tentatives refusées en 401 et 403 \emph{sont} enregistrées, ce qui est
« précieux pour la supervision de sécurité ».
\end{keybox}
